
Empirical calibration anchored the simulation utilizing the Municipal Environment and Natural Resources Office (MENRO) audits, establishing a valid Status Quo baseline of 12.5\% global compliance across the municipality. This objective benchmark was critical in counteracting the heavily inflated self-reported compliance rates of over 50\% found in preliminary surveys \citep{Paigalan2025}. 

To formally evaluate the structural validity and predictive reliability of the developed Agent-Based Model (ABM) against this empirical baseline, a variance-based Global Sensitivity Analysis (GSA) was executed using the Sobol method \citep{Sobol2001}. Unlike local sensitivity analysis—which alters one variable at a time—GSA simultaneously perturbs the entire multi-dimensional input space, allowing for the quantification of both direct effects and non-linear interactions. This computational stress-testing phase involved executing 10,000 parameterized simulation samples, varying the core Theory of Planned Behavior (TPB) parameters and behavioral decay rates.

\begin{table}[htbp]
\caption{Sobol Global Sensitivity Indices ($N=10,000$ Samples)}
\label{tab:sobol_analysis}
\begin{center}
\begin{tabular}{|l|c|c|}
\hline
\textbf{Input Parameter} & \textbf{First-Order ($S_{1}$)} & \textbf{Total-Order ($S_{T}$)} \\
\hline
Attitude Weight ($w_{A}$) & 0.00 & 0.10 \\
Subj. Norm Weight ($w_{SN}$) & 0.06 & 0.21 \\
PBC Weight ($w_{PBC}$) & 0.00 & 0.01 \\
Cost of Effort ($C_{Effort}$) & 0.79 & 1.06 \\
Decay Rate ($\lambda_{dec}$) & 0.01 & 0.00 \\
\hline
\end{tabular}
\end{center}
\end{table}

The extraction of the Sobol indices, summarized in Table \ref{tab:sobol_analysis}, provides a profound mathematical validation of the simulation's underlying psychosocial architecture. The analysis unequivocally reveals that the internal Cost of Effort ($C_{Effort}$) is the overwhelmingly dominant driver of household behavior. Yielding a Total-Order index ($S_{T}$) of 1.06 and accounting for 79\% of the variance in global compliance ($S_{1} \approx 0.79$), this parameter mathematically models the well-documented ``Intention-Action gap'': if the internal friction of washing, sorting, and storing waste is too high, households default to non-compliance regardless of their pro-environmental beliefs. 

Rather than an emergent computational anomaly, this finding serves as a direct validation of the empirical disparities observed across the seven barangays. The baseline behavioral profiles were strictly calibrated to reflect these localized logistical realities. For instance, empirical data indicates a stark inverse correlation between physical friction and compliance: Barangay Poblacion exhibits an extreme effort barrier ($C_{Effort} = 0.75$) due to entrenched mixed-garbage practices, resulting in a baseline compliance of merely 2\%. Conversely, Barangay Babalaya, where waste segregation is culturally integrated and physical effort is low ($C_{Effort} = 0.40$), maintains a near-perfect 100\% compliance rate. This is highly consistent with recent Philippine barangay-level case studies highlighting structural friction as the primary barrier to ecological solid waste management \citep{Yazawa2025}.

Conversely, the intrinsic Attitude weight ($w_{A}$) exhibited near-zero first-order sensitivity ($S_1 = 0.00$), proving that Information, Education, and Communication (IEC) campaigns suffer from strict diminishing returns if implemented without localized structural support. Interestingly, the Subjective Norm Weight ($w_{SN}$) showed moderate interactive effects ($S_T = 0.21$), suggesting that peer pressure and community \textit{Bayanihan} can influence behavior, but only as a secondary catalyst once the primary $C_{Effort}$ bottleneck is addressed.

Consequently, the sensitivity analysis establishes a critical and mathematically rigid constraint for the Mayor Agent. Because the Deep Reinforcement Learning (DRL) agent in this architecture is constrained by heuristic rules rather than independently discovering an unguided strategy from scratch, the GSA results directly inform its action masking and reward structure. To achieve systemic stabilization, the heuristics dictate that the AI cannot rely on passive IEC campaigns to merely boost psychological metrics. Instead, the agent is forced to dynamically deploy tangible, heavy-handed interventions—such as optimized financial incentives or strict penalties—that carry enough mathematical weight to offset the severe $C_{Effort}$ bottlenecks present in critical areas like Barangay Poblacion.

\subsection{The ``Convenience Barrier'' and the Absolute Dominance of Physical Friction}

The most significant revelation of the GSA was the absolute dominance of the Base Cost of Effort ($c_{\text{effort}}$) as the primary driver of system variance. This parameter yielded a First-Order index ($S_1 \approx 0.79$) and a Total-Order index ($S_T \approx 1.06$) that eclipsed all other variables combined. % (Note: Total-order indices marginally exceeding 1.0 are common artifacts of numerical approximation in Saltelli sampling, but fundamentally indicate that nearly 100\% of systemic variance involves this parameter).

These results strongly support the ``Convenience Hypothesis'': in the context of Bacolod, the physical inconvenience of segregating waste—represented by temporal, physical, and cognitive friction—is the ultimate bottleneck. Computationally, this suggests that regardless of a household's pro-environmental beliefs, if the effort threshold exceeds their perceived utility, the behavior will fail \citep{Corcoran2020}. The AI learned that it could not simply ``educate'' away physical friction. Instead, it had to fundamentally alter the individual utility calculus by using Enforcement to make non-compliance costlier and Incentives to make compliance more rewarding.

\subsection{The ``Value-Action Gap'' and the Intention-Action Trap}

A critical corollary to the dominance of effort was the near-zero direct sensitivity of the Attitude ($w_a$) and Perceived Behavioral Control ($w_{pbc}$) parameters. Both exhibited First-Order indices ($S_1$) of approximately 0.01. While traditional municipal strategies focus heavily on IEC campaigns, the model demonstrates a computational realization of the ``Value-Action Gap'' \citep{Liao2024, Zhao2022}.

Interestingly, while Attitude had a negligible direct effect, its Total-Order index ($S_T \approx 0.09$) indicates that its influence is almost entirely driven by non-linear interactions with other variables (likely Cost of Effort and Social Norms). Nevertheless, the GSA confirms that passive awareness campaigns are mathematically insufficient to drive compliance if the systemic effort ($c_{\text{effort}}$) remains unaddressed. Consequently, the HuDRL agent's strategic pivot away from IEC in its early phases was a calculated recognition that the most effective lever for immediate behavioral shifts lies in structural motivators rather than intrinsic appeals.

\subsection{Sensitivity Analysis of Social Tipping Points}
To address the concern that the efficacy of the proposed policy might be contingent upon specific "tuned" behavioral parameters, a hardcoded robustness test was performed. The two primary behavioral thresholds—the compliance floor (\(T_{low}\)) and the social norm tipping point (\(T_{high}\))—were varied across a matrix of nine scenarios (Table \ref{tab:robustness_results}). These scenarios ranged from optimistic social environments (\(T_{high} = 0.60\)) to extreme stubbornness (\(T_{high} = 0.85\)).

The results demonstrate a high degree of structural resilience in the HuDRL-driven Sequential Saturation strategy. Despite a 25\% increase in the social resistance threshold (\(T_{high}\) shifting from 0.60 to 0.85), the final municipality-wide compliance remained remarkably stable, fluctuating only between 89.83\% and 91.84\%.

This stability suggests that the success of the policy is not a finding "tuned" to a specific threshold, but rather a result of the adaptive nature of the Deep Reinforcement Learning agent. The Mayor Agent, utilizing the Proximal Policy Optimization (PPO) algorithm, monitors real-time compliance states and adaptively adjusts the intensity and duration of interventions in each barangay. In "stubborn" scenarios with higher tipping points, the RL agent compensates by sustaining resource allocation until the specific social threshold of the community is surpassed, effectively "locking in" the habit shield regardless of the initial social friction. This robustness confirms that the Sequential Saturation strategy is a structurally valid finding that is resilient to variations in underlying behavioral assumptions.

\begin{table}[htbp]
\centering
\caption{Hardcoded Robustness Test Results for Tipping Point Variance}
\label{tab:robustness_results}
\begin{tabular}{@{}cccc@{}}
\toprule
\textbf{Scenario} & \textbf{\(T_{low}\) (Inertia)} & \textbf{\(T_{high}\) (Shield)} & \textbf{Final Compliance (\(C_f\))} \\ \midrule
1                 & 0.20                           & 0.60                            & 91.00\%                            \\
2                 & 0.25                           & 0.70                            & 91.80\%                            \\
3                 & 0.30                           & 0.80                            & 91.22\%                            \\
4                 & 0.35                           & 0.60                            & 90.68\%                            \\
5                 & 0.40                           & 0.70                            & 90.45\%                            \\
6                 & 0.45                           & 0.80                            & 91.00\%                            \\
7                 & 0.50                           & 0.70                            & 91.84\%                            \\
8                 & 0.55                           & 0.80                            & 89.83\%                            \\
9                 & 0.60                           & 0.85                            & 91.35\%                            \\ \bottomrule
\end{tabular}
\end{table}

\subsection{Non-Linear Interactivity and Habit Decay}

The GSA demonstrated significant non-linear interactivity, particularly concerning the Cost of Effort. The substantial delta between its First-Order and Total-Order indices suggests that a household's threshold for effort does not exist in a vacuum; it is heavily entangled with the other psychosocial metrics and the localized environment. 

Conversely, the Habit Decay rate ($\delta$) exhibited near-zero influence on overall variance ($S_1 \approx 0.00, S_T \approx 0.01$). This provides a critical sociological insight: citizens are not failing to segregate because they ``forget'' the rules over time; they are actively and rationally choosing not to comply because the physical effort remains consistently too high. 

\subsection{Limitations and the Informal Economy}

While the GSA clearly identifies Cost of Effort as the dominant variable, the ``Net Cost'' ($C_{Net}$) of non-compliance in the real world is often distorted by the informal economy. As established in the key informant interviews, practices such as localized bribery or ``tipping'' collection personnel to bypass the ``no segregation, no collection'' rule can artificially lower the perceived cost of non-compliance, effectively functioning as an informal reduction in $c_{\text{effort}}$.

If a nominal bribe effectively bypasses the official penal structure, it undermines the utility calculus established by the DRL strategy. This implies that for the HuDRL policies to be effective in reality, the LGU must couple its resource allocation with strict auditing of collection crews and the elimination of informal bypasses. Without ensuring that the simulation's enforcement parameters match on-the-ground reality, the dominant ``Convenience Barrier'' identified by the Sobol analysis will remain an insurmountable obstacle to municipal sustainability.